\clearpage\pagestyle{otMain}\setlength{\otBlocEspacement}{2.5mm}\fontsize{11}{15}\selectfont\corrigesChapitre{5}{Égalités, programmes et suites}
\begin{blocLivre}{cor-5-1}\begin{corrige}{1}{Égalités à trous}\demarche{Calculer d’abord le côté connu.}
$6\times8=48$ puis $178-48=\res{130}$. Ensuite $25\times4=\res{100}$. Enfin $630:9=\res{70}$.
\end{corrige}\end{blocLivre}
\begin{blocLivre}{cor-5-2}\begin{corrige}{2}{Exécuter le programme}\demarche{Suivre les étapes dans l’ordre.}
Avec 6 : $6+2=8$, $8\times4=32$, $32-3=\res{29}$. Avec 10 : $12\times4=48$, $48-3=\res{45}$.
\label{fig-cor-5-2}\figureGeo[corrige]{prog-2}
\end{corrige}\end{blocLivre}
\begin{blocLivre}{cor-5-3}\begin{corrige}{3}{Remonter le programme}\demarche{Défaire les étapes en sens inverse.}
On remonte : $48:2=24$ ; $24-6=18$ ; $18:3=\res6$. Contrôle : $(6\times3+6)\times2=48$.
\end{corrige}\end{blocLivre}
\begin{blocLivre}{cor-5-4}\begin{corrige}{4}{Deux programmes, deux sorties}\demarche{Tester chaque programme séparément.}
Premier programme : $\res{32}$ puis $\res{48}$. Second : $6\times4+2=\res{26}$ puis $10\times4+2=\res{42}$. L’ordre des étapes change la sortie.
\end{corrige}\end{blocLivre}
\begin{blocLivre}{cor-5-5}\begin{corrige}{5}{La règle de la suite}\demarche{Comparer les écarts, puis les rapports.}
Première suite : on ajoute 8 ; suivent $\res{39}$ et $\res{47}$. Seconde : chaque terme double ; suivent $\res{80}$ et $\res{160}$.
\end{corrige}\end{blocLivre}
\begin{blocLivre}{cor-5-6}\begin{corrige}{6}{Double plus un}\demarche{Tester la règle sur chaque passage.}
$7\times2+1=15$ ; $15\times2+1=31$ ; $31\times2+1=63$ : la règle convient. Ensuite : $63\times2+1=\res{127}$, puis $127\times2+1=\res{255}$.
\end{corrige}\end{blocLivre}
\begin{blocLivre}{cor-5-7}\begin{corrige}{7}{Le nombre caché}\demarche{Remonter les opérations une à une.}
Avant de retirer 7, le produit valait $41+7=48$ ; avant de multiplier, le nombre était $48:6=\res8$. Contrôle : $8\times6-7=41$.
\end{corrige}\end{blocLivre}
\begin{blocLivre}{cor-5-8}\begin{corrige}{8}{Produire un programme}\demarche{Partir de la sortie et chercher un produit proche.}
Par exemple « multiplier par 5, ajouter 3 » : $4\times5+3=\res{23}$. Autre programme en deux étapes : « multiplier par 6, retirer 1 » : $4\times6-1=\res{23}$. Plusieurs chemins mènent à la même sortie.
\end{corrige}\end{blocLivre}
\begin{blocLivre}{cor-5-9}\begin{corrige}{9}{La suite qui zigzague}\demarche{Observer les écarts un pas sur deux.}
La règle alterne $+9$ puis $-3$. Après 71 : $71-3=\res{68}$, puis $68+9=\res{77}$, puis $77-3=\res{74}$.
\end{corrige}\end{blocLivre}
\begin{blocLivre}{cor-5-10}\begin{corrige}{10}{Des symboles cohérents}\demarche{Résoudre d’abord l’égalité à un seul symbole.}
Trois triangles font 27 : $\triangle=\res9$. Alors $9+\bullet=17$ donne $\bullet=\res8$.
\end{corrige}\end{blocLivre}
\begin{blocLivre}{cor-5-11}\begin{corrige}{11}{Le motif qui grandit}\demarche{Compter les ajouts depuis l’étape 1.}
Étape 10 : neuf ajouts après la première : $5+9\times4=\res{41}$ carreaux. Pour 33 carreaux : $33-5=28$, soit $28:4=7$ ajouts : étape $\res8$. Contrôle : $5+7\times4=33$.
\label{fig-cor-5-11}\figureGeo[corrige]{motifs-11}
\end{corrige}\end{blocLivre}
\begin{blocLivre}{cor-5-12}\begin{corrige}{12}{Les tee-shirts, suite}\demarche{Utiliser la formule dans les deux sens.}
Pour 7 : $7\times12+5=\res{89}$ euros. Pour 125 euros : les tee-shirts coûtent $125-5=120$ euros, donc $N=120:12=\res{10}$.
\end{corrige}\end{blocLivre}
\begin{blocLivre}{cor-5-13}\begin{corrige}{13}{Le tableau de la machine}\demarche{Tester une règle sur toutes les lignes du tableau.}
D’une entrée à l’autre, la sortie monte de 5 par unité d’entrée : essai « multiplier par 5, ajouter 2 » : $3\times5+2=17$ ; $5\times5+2=27$ ; $8\times5+2=42$ : la règle convient. Pour 20 : $20\times5+2=\res{102}$.
\end{corrige}\end{blocLivre}
\begin{blocLivre}{cor-5-14}\begin{corrige}{14}{Vrai ou faux ?}\demarche{Chercher deux règles compatibles avec le début.}
« Ajouter 2 » donne $2 ;\ 4 ;\ \res6 ;\ 8$. « Doubler » donne $2 ;\ 4 ;\ \res8 ;\ 16$. Deux termes ne suffisent pas : l’affirmation est $\res{\text{fausse}}$. Une règle se vérifie sur plusieurs passages.
\end{corrige}\end{blocLivre}
\begin{blocLivre}{cor-5-15}\begin{corrige}{15}{Le défi du professeur}\demarche{Chercher un nombre multiplié par lui-même.}
Le nombre multiplié par lui-même donne 36 : $6\times6=36$. Le candidat 6 est pair : $\square=\res6$. C’est la seule solution entière positive paire.
\end{corrige}\end{blocLivre}
